\songtitle{Line of Fire}

\setin{2005-08-31}
\tag*{2026}\tag{salsa}\tag{bolero}\tag{ricochets}\tag{miguel-morocho}\tag{john-carlson}
% -> ricochets.tex
\begin{notes}
Lyrics by Claude based on \emph{Ricochets} by Carlos Thompson -- chapter 13 (\emph{The Glades}): in Miguel Morocho's voice, in his fifth-floor office in Cali, the plastic bag on the desk (a passport, credit cards, a Motorola Razr), stolen by petty muggers and recovered by his people, rum against Carlson's Coke, ``She put herself in the line of fire'', and lunch. Retrieved from rataflechera's Suno page (V6, 24 November 2026).
\end{notes}
\begin{style}
Cali salsa bolero, velvet piano tumbao, muted trumpet lead, soft congas and bongó, güiro, warm upright bass, opening nylon-guitar bolero rising into a restrained coro salsa; hushed close-mic male baritone with soft harmony pads beneath the chorus, smooth and velvety, courteous and unhurried; 92 BPM, slow son montuno feel in A minor.
\end{style}
\begin{lyrics}
\songpart{Intro}
\annotation{nylon guitar bolero, a ceiling fan, ice dropped into a glass}

\songpart{Verse 1}
\annotation{male:}\\
Siéntese, John, the fifth floor is cooler,\\
the river's out there doing what it does.\\
A Viejo de Caldas for me, a Coke for you,\\
I remember how you take it, como era.\\
On the desk, a plastic bag, nothing fancy,\\
I'd rather hand it over than explain.\\
A passport. Credit cards. A Motorola Razr.\\
Go on. Take your time. I'll pour again.

\songpart{Pre-Coro}
\annotation{male:}\\
Little hands took them off the concrete,\\
and the world kept walking with them for a while.\\
One of them got as far as Panama,\\
someone tried to wear her face and smile.\\
My people brought them home.

\songpart{Coro}
\annotation{coro:}\\
En la línea de fuego
\annotation{male:}\\
she put herself, John, nobody pushed her there,
\annotation{coro:}\\
en la línea de fuego
\annotation{male:}\\
she knew the price and she paid it fair.
\annotation{coro:}\\
Ay, la niña, la niña
\annotation{male:}\\
the daughter, no one anticipated,
\annotation{coro:}\\
nadie la vio venir
\annotation{male:}\\
not me, not them, not you.

\songpart{Verse 2}
\annotation{male:}\\
You don't open it. That's all right, hermano,\\
some things are heavy even through the plastic.\\
Her picture's looking up at the ceiling fan,\\
the phone is dead, I didn't want to charge it.\\
I won't tell you who and I won't tell you why,\\
you know the trade, you know the way it's paid.\\
I only tell you what I'd want to hear\\
if it were mine, if I had stayed.

\songpart{Pre-Coro}
\annotation{male:}\\
Little hands took them off the concrete,\\
and the world kept walking with them for a while.\\
One of them got as far as Panama,\\
someone tried to wear her face and smile.\\
My people brought them home.

\songpart{Coro}
\annotation{coro:}\\
En la línea de fuego
\annotation{male:}\\
she put herself, John, nobody pushed her there,
\annotation{coro:}\\
en la línea de fuego
\annotation{male:}\\
she knew the price and she paid it fair.
\annotation{coro:}\\
Ay, la niña, la niña
\annotation{male:}\\
the daughter, no one anticipated,
\annotation{coro:}\\
nadie la vio venir
\annotation{male:}\\
not me, not them, not you.

\songpart{Bridge}
\annotation{the band drops to piano and congas, the trumpet muted and far}
\annotation{male:}\\
Five minutes, John. I'll give you five.\\
I'll watch the river, I won't watch your hands.\\
The rum is warm, your Coke has lost its fizz,\\
and nothing on this desk will understand.

\annotation{Mambo}
\annotation{trumpet line, restrained, piano montuno}
\annotation{coro:}\\
En la línea de fuego, en la línea de fuego,
\annotation{male:}\\
nobody pushed her, nobody pushed her there.
\annotation{coro:}\\
En la línea de fuego, nadie la vio venir,
\annotation{male:}\\
not me, not them, not you.

\songpart{Outro}
\annotation{back to the bolero guitar, the fan, a chair pushed back}
\annotation{male:}\\
Bueno.\\
Let's get some lunch.
\annotation{End}
\end{lyrics}

